\originalchapter{局域矩、Friedel相移与Kondo效应}
\label{ch:kondo}
\source{8.512第20--22讲，含两个文件名不同的第22讲同文版本。}
\originalsection{散射相移与粒子数}
球对称杂质势使每个角动量通道的远区径向波为$\sin(kr-l\pi/2+\delta_l(k))$。在半径$R$盒子中边界量子化$kR-l\pi/2+\delta_l(k)=n\pi$，相移使低于$k_F$的态数变化$\Delta N_l=\delta_l(k_F)/\pi$，包含相应束缚态的计数后成立。每个$l$有$2l+1$角简并及自旋，故$\Delta N=\sum_{l\sigma}(2l+1)\delta_{l\sigma}(E_F)/\pi$。这是Friedel和规则，连接局域屏蔽电荷与散射相移。

例如局域$s$波共振在费米能处$\delta=\pi/2$，每自旋贡献半个电子的相移计数。谱峰的宽度由杂质--传导带杂化决定，不能把半个计数当作杂质可单独处于“半个电子”的本征态。
\originalsection{Anderson模型与局域矩条件}
Anderson杂质模型为
\begin{equation}
H=\sum_{k\sigma}\epsilon_kc_{k\sigma}^\dagger c_{k\sigma}+\epsilon_d\sum_\sigma n_{d\sigma}+Un_{d\uparrow}n_{d\downarrow}+\sum_{k\sigma}(V_kc_{k\sigma}^\dagger d_\sigma+\mathrm{h.c.}).
\end{equation}
无杂化时空、单、双占据能为$0,\epsilon_d,2\epsilon_d+U$。以$E_F=0$计，$\epsilon_d<0<\epsilon_d+U$使单占据最低；若两种电荷代价远大于杂化宽度$\Gamma=\pi D_\uparrow|V|^2$，可形成局域矩区。$U=0$时$G_d^R(\omega)=1/(\hbar\omega-\epsilon_d+\ii\Gamma)$，谱函数Lorentz形，$n_{d\sigma}=1/2-\pi^{-1}\arctan(\epsilon_d/\Gamma)$。

静态Hartree近似用$\epsilon_{d\sigma}=\epsilon_d+Un_{d,-\sigma}$，可预测磁性分裂，但它把量子自旋翻转抑制成静态场，无法得到低温Kondo单重态。判定局域矩形成的中间温区与最终零温屏蔽要分开。
\originalsection{Schrieffer--Wolff与对数流}
在局域矩区，把空占据和双占据作为虚中间态，用二阶有效哈密顿量得到$H_K=J_K\vct S_d\cdot\vct s(0)+W\rho(0)$，
\begin{equation}
J_K=2|V|^2\left(\frac1{|\epsilon_d|}+\frac1{\epsilon_d+U}\right)>0,\quad
W=\frac{|V|^2}{2}\left(\frac1{|\epsilon_d|}-\frac1{\epsilon_d+U}\right),
\end{equation}
这里$\vct s=\frac12c^\dagger\vct\sigma c$，$W$的系数随局域场归一而变，但$J_K$反铁磁性及粒子--空穴对称时$W=0$不变。虚跃迁把两个自旋交换，解释电荷虽冻结而自旋仍与金属耦合。

弱耦合散射的粒子和空穴虚过程给$J_K^2D_\uparrow\log(D/E)$修正。令$l=\log(D_0/D)$，常用约定$j=D_\uparrow J_K$满足$\dd j/\dd l=2j^2+O(j^3)$，解$j(l)=j_0/(1-2j_0l)$。正耦合随降低能量增长，尺度$\kb T_K\sim D_0\exp[-1/(2j_0)]$。当$j\sim1$，该微扰方程已失效；形式发散不是物理量真的无穷大，而是强耦合新描述开始的信号。
\originalsection{单重态与低温共振}
单通道、自旋$1/2$、反铁磁耦合的强耦合极限中，杂质与局域传导自旋形成单重态，能量$-3J_K/4$；三重态为$J_K/4$。低温保留的传导粒子散射有$\pi/2$相移，并有弱的残余相互作用。零温杂质谱在$E_F$附近产生宽约$\kb T_K$的Kondo共振，它不是裸$\epsilon_d$电荷共振的平移。

金属电阻在中间温区有随降温上升的$\log T$修正，低温饱和并有$T^2$修正。磁化率趋有限常数而非继续Curie发散，杂质比热$C_{
m imp}\propto T$。单通道完全屏蔽问题可用局域Fermi液体解释；过屏蔽多通道、自旋大于可用通道能力或伪隙态密度的问题可能不同，本章结论不覆盖这些模型。
